\section{General k-injection details}
\label{app:k-injection-details}

In this section, we give the stabilizer-cleaning construction used in \Cref{sec:k-injection} and explain why it preserves the encoded logical information.

\subsection{Logical representatives reduction}
We describe a recursive algorithm for transforming the logical representatives to a set of $ k $ physical qubits.
Let
\begin{equation}
\mathcal{L} := \{ \bar{X}_i, \bar{Z}_i \}_{i=1}^{k}
\label{eq:tomas-current-logical-basis}
\end{equation}
be the current logical representatives, and let $\mathcal{S}$ be the current stabilizer generator
set.

The recursive reduction is as follows. Given a non-trivial stabilizer group, 
% that is one where there must exist at least one stabilizer of weight greater than 1, 
do the following:
\begin{enumerate}
    \item Pick a stabilizer $s \in \mathcal{S}$ with $\operatorname{wt}(s) > 1$.
    \item Measure a single-qubit Pauli $P_j$ that anticommutes with $s$ on a qubit $j \in \operatorname{supp}(s)$, i.e.,
    \begin{equation}
    \{ P_j, s \} = 0;
    \label{eq:tomas-step-anticommute}
    \end{equation} denote the measurement outcome by $m \in \{ \pm 1 \}$.
    \item Replace $s \in \mathcal{S}$ by the measured $P_j$.
    \item Update any remaining stabilizers and logical representatives that
    anticommute with $P_j$ by multiplying them by the replaced stabilizer $s$.
\end{enumerate}

To write this explicitly, let
\begin{equation}
\mathcal{S} \setminus \{ s \} = \{ g_a \}_{a=1}^{n-k-1} .
\label{eq:tomas-step-other-stabilizers}
\end{equation}
The updated generator set is
\begin{equation}
\mathcal{S}' := \{ m P_j \} \cup \{ g_a' \}_{a=1}^{n-k-1},
\label{eq:tomas-step-stabilizer-set}
\end{equation}
where
\begin{equation}
g_a' =
\begin{cases}
g_a, & [g_a, P_j] = 0, \\[3pt]
g_a s, & \{ g_a, P_j \} = 0 .
\end{cases}
\label{eq:tomas-step-stabilizer-update}
\end{equation}
Similarly, each logical representative $L \in \mathcal{L}$ is updated to
$L' \in \mathcal{L}'$ according to
\begin{equation}
L' =
\begin{cases}
L, & [L, P_j] = 0, \\[3pt]
L s, & \{ L, P_j \} = 0 .
\end{cases}
\label{eq:tomas-step-logical-update}
\end{equation}
We denote the updated logical representatives by
\begin{equation}
\mathcal{L}' := \{ \bar{X}_i', \bar{Z}_i' \}_{i=1}^{k}.
\label{eq:tomas-step-updated-logical-basis}
\end{equation}

Each reduction ``peels off'' a physical qubit $j$ from the original logical representatives $\mathcal{L}$, since qubit $j$ is fixed by a weight-$1$ Pauli $P_j$.
We repeat the above reduction step
until all stabilizers are weight-$1$, and obtain
\begin{equation}
\mathcal{S} \to \mathcal{S}', \quad
\mathcal{L} \to \mathcal{L}',
\end{equation}
where $ \mathcal{S}' $ contains only weight-$1$ Pauli operators,
and $ \mathcal{L}' $ is the new set of logical representatives that only live on a set of $ k $ physical qubits. 
We call these $ k $ physical qubits the \emph{carrier qubits} and denote them by $\mathcal{Q}_L = \set{Q_i}_{i=1}^k$, and call the remaining $ n-k $ physical qubits supporting single-qubit Paulis in $ \mathcal{S}' $ the \emph{peeled qubits}, $\mathcal{Q}_P$.

Since $ \mathcal{L}' $ is a valid conjugate logical representative,
there exists a Clifford $C$ such that it maps the standard computational basis on the $ k $ carrier qubits to the new logical representatives $\mathcal{L}'$:
\begin{equation}
C X_{Q_i} C^\dagger = \bar{X}_i',
\qquad
C Z_{Q_i} C^\dagger = \bar{Z}_i'.
\end{equation}
And thus, $C Y_{Q_i} C^\dagger = \bar{Y}_i'$.

Therefore, to prepare the $+1$ eigenstate of $ \bar{M}_i' $ for all $ i \in [k] $, we first prepare the carrier qubits $\set{Q_i}$ in the $+1$ eigenstates of $ \{M_{Q_i}\} $, then apply the Clifford circuit $C$ to the carrier qubits:
\begin{equation}
\begin{aligned}
    \bar{M}_i' &= \cos\theta_i\, \bar{X}_i' + \sin\theta_i\, \bar{Y}_i' \\
    &= C \bigl(\cos\theta_i\, X_{Q_i} + \sin\theta_i\, Y_{Q_i}\bigr) C^\dagger \\
    &= C \;{M}_{Q_i}\; C^\dagger .
\end{aligned}
\end{equation}

Although the construction above applies generally to arbitrary CSS codes and can inject all \(k\) logical qubits simultaneously, it has several caveats.
One limitation is that the state-preparation Clifford circuit \(C\) is typically nonlocal and deep, and therefore not itself fault tolerant. 
This issue may be mitigated by implementing or compiling \(C\) on the logical level instead of the physical level. 
For instance, on an extractor architecture~\cite{he_extractors_2025} implementing Pauli-based computation~\cite{Bravyi2016}, a logical Clifford circuit can be implemented by changing the Pauli basis of subsequent logical Pauli-rotations. 
In other words, instead of implementing \(C\), we simply change subsequent gates from \(e^{i\theta P}\) to \(e^{i\theta C^\dagger PC}\). 
In a model where every Pauli rotation has roughly the same implementation cost~\cite{he_extractors_2025}, this Clifford rotation is essentially free.

\subsection{Logical information preservation}

During the reduction step, the update rule in
\Cref{eq:tomas-step-stabilizer-update,eq:tomas-step-logical-update}
preserves the logical subsystem for three reasons.

First, the measurement of $P_j$ reveals no logical information. Let
$\ket{\phi}$ be any state in the codespace of $S$, so in particular
$s \ket{\phi} = \ket{\phi}$. Using \Cref{eq:tomas-step-anticommute},
\begin{equation}
\bra{\phi} P_j \ket{\phi}
=
\bra{\phi} s P_j s \ket{\phi}
=
- \bra{\phi} P_j \ket{\phi}
= 0 .
\label{eq:tomas-step-random-outcome}
\end{equation}
Hence the two outcomes $m = \pm 1$ occur with equal probability $1/2$,
independently of the encoded logical state. The measurement therefore creates a
new local stabilizer sign, but it does not read out any logical observable.

Second, the number of encoded qubits is unchanged. The original generator $s \in \mathcal{S}$ is
removed and replaced by the new independent generator $m P_j$. Thus
$\mathcal{S}'$ again contains $n-k$ independent commuting generators, so the
updated code still encodes exactly $k$ logical qubits.

Third, the logical representatives are changed only by multiplication with the original
stabilizer $s$. If $L' = L s$, then on the original codespace,
\begin{equation}
L' \ket{\phi} = L s \ket{\phi} = L \ket{\phi} .
\label{eq:tomas-step-same-logical-action}
\end{equation}
Thus \Cref{eq:tomas-step-logical-update} changes only the
\emph{representatives} of the logical Paulis, not the logical Pauli classes
themselves.
And thus, 
\begin{equation}
[\bar{M}_i', \proj_{\mathcal{S}}] = 0 \quad \forall i \in [k].
\end{equation}
In another word, if the state is prepared in the simultaneous $ +1 $ eigenstate of $ \bar{M}_i' $ for all $ i \in [k] $,
then after the SE of $\mathcal{S}$, the state is the simultaneous $ +1 $ eigenstate of $ \bar{M}_i $ for all $ i \in [k] $.
